% Figure: global error at t=5 versus step size, log-log, for forward Euler
% and RK4 on the test problem ydot = -y, y(0)=1. Slopes 1 and 4, with the
% RK4 roundoff floor at small h.
\begin{figure}[htbp]
\centering
\begin{tikzpicture}
\begin{axis}[
  width=12cm, height=8cm,
  xlabel={step size $h$}, ylabel={absolute error at $t=5$},
  xmode=log, ymode=log,
  xmin=3e-4, xmax=1, ymin=1e-16, ymax=1e-1,
  axis lines=left,
  legend pos=south east,
  legend cell align=left,
  grid=major, grid style={enggray!20},
  tick label style={font=\scriptsize},
  label style={font=\small},
  legend style={font=\scriptsize},
]
% Forward Euler measured points (from the chapter worked example / exercise table)
\addplot[engred, very thick, mark=*, mark options={fill=engred}] coordinates {
  (0.5, 5.76e-3)
  (0.1, 1.59e-3)
  (0.01, 1.69e-4)
  (0.001, 1.68e-5)
};
\addlegendentry{forward Euler (measured)}
% Slope-1 reference line through (0.01, 1.69e-4): err = 0.0169 * h
\addplot[enggray, dashed, thin, domain=0.0005:0.7] {0.0169*x};
\addlegendentry{slope 1 reference}
% RK4 measured points
\addplot[engblue, very thick, mark=square*, mark options={fill=engblue}] coordinates {
  (0.5, 5.1e-5)
  (0.1, 1.39e-7)
  (0.01, 1.41e-11)
  (0.001, 1.0e-13)
};
\addlegendentry{RK4 (measured)}
% Slope-4 reference through (0.1, 1.39e-7): err = 1.39e-3 * h^4
\addplot[enggray, dotted, thin, domain=0.004:0.7] {1.39e-3*x^4};
\addlegendentry{slope 4 reference}
% roundoff floor annotation
\node[engblue, font=\scriptsize, anchor=south] at (axis cs:0.0015,3e-13) {roundoff floor};
\end{axis}
\end{tikzpicture}
\caption[Global error versus step size for Euler and RK4]{Absolute
error at $t=5$ for the test problem $\dot{y}=-y$, $y(0)=1$ (exact
$e^{-5}\approx 6.74\times 10^{-3}$), as a function of step size on
log-log axes. Forward Euler (red circles) falls on a line of slope
$1$: halving $h$ halves the error, the signature of a first-order
method. RK4 (blue squares) falls on a line of slope $4$: halving $h$
cuts the error by a factor of $16$, until $h\approx 10^{-3}$ where
the truncation error drops below the accumulated floating-point
roundoff and the curve flattens at the precision floor
$\sim N\epsilon\approx 10^{-12}$ to $10^{-13}$. The figure shows the
two facts an engineer needs about a numerical method at once: its
order of accuracy (the slope) and the arithmetic floor below which
no step refinement helps. Computed from the closed-form
amplification factors $y_{N}=(1-h)^{5/h}$ for Euler and
$y_{N}=R(h)^{5/h}$ for RK4, where $R(h)$ is the truncated
exponential.}
\label{fig:vol02:ch07:euler-rk4-error}
\end{figure}
